%% LeafletFA -- Bioinformatics (OUP) Original Paper, companion to SpliceVI (submitted together).
%% Skeleton created 2026-11-01 from BIOINFORMATICS_PLAN.md. Source material: GenomeBiology_Submission2026/ (never submitted; kept).

%% Layout: same class options as the SpliceVI paper (SpliceVI_paper_overleaf/bioinformatics_main.tex):
%% traditional + large = 178 mm text width, two 86 mm columns; namedate = author-year citations.
\documentclass[webpdf,traditional,large,namedate]{oup-authoring-template}

\usepackage[utf8]{inputenc}
\usepackage[T1]{fontenc}
\usepackage{url}
\usepackage{xurl}
\usepackage{booktabs}
\usepackage{float}
\usepackage{graphicx}
\usepackage{amsfonts}
\usepackage{amsmath}
\usepackage{microtype}
\usepackage{xcolor}
\usepackage{xr}
\externaldocument[S-]{leafletfa_supplement} % supplement labels: \ref{S-<label>}

\newcommand{\societylogo}{}
\let\cite\citep

% Result numbers come from numbers.tex (generated from paper/ledger/results_ledger.tsv; do not hand-edit).
% Red fallbacks show until numbers.tex defines the macro.
\InputIfFileExists{numbers}{}{}
\providecommand{\DuplicateAllOnOne}{\textcolor{red}{68\%}}
\providecommand{\DuplicateMolecules}{\textcolor{red}{1.7--3.1}}

% Drafting markup (red in the PDF); remove before submission
\newcommand{\unsure}[1]{\textcolor{red}{#1}}
\newcommand{\todocite}[1]{\textcolor{red}{[cite: #1]}}
\newcommand{\todo}[1]{\textcolor{red}{[TODO: #1]}}

\begin{document}

\journaltitle{Bioinformatics}
\DOI{DOI added during production}
\copyrightyear{2026}
\pubyear{2026}
\vol{XX}
\issue{x}
\access{Published: Date added during production}
\appnotes{Original Paper}

\firstpage{1}

%\subtitle{Subject Section} % TODO: Bioinformatics subject section (Gene expression, or Sequence analysis)

% Working title (Karin, 2026-11-01); final after check 10 (amendment 6). Original Paper; Application Note as fallback.
\title[LeafletFA]{A cell-type-resolved atlas of alternative splicing in mouse full-length single-cell data with LeafletFA}

\author[1]{\todo{authors and order}}
\address[1]{\todo{affiliations}}
\corresp[$\ast$]{Corresponding author. \todo{email}}

\received{Date}{0}{Year}
\revised{Date}{0}{Year}
\accepted{Date}{0}{Year}

\abstract{\mdseries
\textbf{Motivation:} \todo{Single-cell splicing is measured as a few duplicated reads per event; per-cell methods need an event definition that works across cells and a likelihood that does not treat duplicate reads as independent evidence.}\\
\textbf{Results:} \todo{ATSEmapper and the ATSE catalog (\CatalogATSEsMouse{} events, \CatalogCellsMouse{} mouse cells); LeafletFA, a Dirichlet-multinomial factor model; its per-cell representation reads junction usage, fits held-out reads better than fair nulls, is stable across seeds, and recovers the generating parameters at full-length depth; programs that recur across seeds are the broad lineages.}\\
\textbf{Availability and Implementation:} \todo{one package (ATSEmapper + LeafletFA), GitHub, PyPI, Zenodo DOI for the catalog and fits}\\
\textbf{Contact:} \todo{email}\\
\textbf{Supplementary Information:} Supplementary data are available at \textit{Bioinformatics} online.}

\maketitle

\section{Introduction}
% P1 why single-cell splicing; P2 what the data are (few duplicated reads per event; detection follows expression, cite SpliceVI2026 for the measurement);
% P3 what exists (LeafCutter pseudobulk, scQuint, SpliceVI as companion); P4 what this paper provides (catalog + model + validation).
\todo{Introduction, four paragraphs (plan section 1).}

\section{Methods}
\subsection{Data}
% TMS and Allen Brain Smart-seq2; one description shared in substance with the companion paper, written fresh here (plan section 3).
\todo{Data.}

\subsection{ATSEmapper and the ATSE catalog}
% Draft 2026-11-01 from DEFINITIONS.md (file:line there); shared with the SpliceVI paper (same wording of the definition,
% each paper's own sentences around it). Numbers via numbers.tex once in the ledger.
Splice junctions were extracted from each cell's alignments with regtools \cite{regtools} and pooled across cells; junctions
with introns of 50--500,000 bp, canonical splice motifs (GT-AG, GC-AG, AT-AC) and at least 100 reads in at least 10 cells
were kept. Each junction was compared with a long-read-supported annotation \cite{Schertzer2025-lw}: an end is annotated
if it lies within 1 bp of an annotated exon boundary on the same strand, and a junction is assigned the genes of the
transcripts it matches. Junctions needed at least one annotated end and a single gene; junctions matching more than one
gene were removed. An alternative transcript structure event (ATSE) is a connected component of the junctions of a gene
in which two junctions are linked if they share a donor or an acceptor site; a junction was removed from its ATSE if it
carried less than 1\% of the reads at either of its splice sites, and an ATSE needs at least two junctions. Junctions that
share no splice site with another junction (singletons) are not alternative within this definition and were dropped,
including constitutive junctions of highly expressed genes. After cell and junction filtering (Supplementary Methods), ATSEs
were re-derived on the surviving junctions; this dropped a further 9,065 singletons, including alternative junctions whose
partners failed the filter of at least 75 cells. A mutually exclusive exon pair forms one ATSE only if another junction links its two
halves; otherwise it appears as two ATSEs, an alternative-3$'$ and an alternative-5$'$ pair (Supplementary
Table~\ref{S-tab:eventtypes}).

\subsection{Model}\label{sec:model}
% Dirichlet-multinomial likelihood, within-ATSE simplex psi, per-ATSE concentration floored at 2 (a modeling choice; Supplementary Note 1),
% Dirichlet prior on phi with concentration pinned at K, alpha_pi ~ Gamma(1,1); stochastic variational inference; frozen-global E-step for new cells.
\todo{Model (from methods\_CURRENT.tex, rewritten to Bioinformatics length).}

\subsection{Ratio-redrawing test}
% Self-contained (Karin, 2026-11-01); the SpliceVI companion paper applies the same test across methods. Settings from
% LeafletFA scripts/splicevi_baseline/knn_check.py and the shared test inputs (SpliceVI STAGE5 §170). \unsure = to verify.
Because a cell detects an event mostly when its gene is expressed, a splicing representation can separate cell types from
which events are detected rather than from how their junctions are used. To measure what the representation reads, we
built a null version of the held-out test cells that keeps, for every cell, which events have reads and how many reads
each has, but removes cell-type information from the junction ratios: the reads of each observed (cell, event) pair were
redrawn from a multinomial with the cell's own read total and the event's population junction ratios
\unsure{(estimated on the training cells)}. A second null set the reads to the read total times the population ratios
exactly, without sampling noise. LeafletFA embedded the real and the null test cells with the same frozen fit (E-step,
Section~\ref{sec:model}), and each cell was represented by its centered log program weights (all $K$ programs; weights
clipped at $10^{-6}$; standardized). Cell-type accuracy is that of a $k$-nearest-neighbor classifier ($k = 15$) with labels
of tissue and broad cell type, in four folds grouped by mouse with a fixed fold assignment, averaged over folds. The drop
in accuracy from the real cells to the redrawn null measures how much of the representation's cell-type structure depends
on junction usage rather than on detection; a drop of at least 0.10 was pre-specified as reading junction usage
\cite{SpliceVI2026}. The companion paper \cite{SpliceVI2026} applies this test to SpliceVI and other splicing models, using
LeafletFA's programs used by cells (mean training weight $> 0.01$); here we report all $K$ programs.

\subsection{Validation design}
% Pre-registered checks (LeafletFA_validation/VALIDATION_PREREG.md): held-out likelihood (event split), kNN with ratios redrawn,
% seed stability, simulation recovery, consensus programs.
\todo{Validation design.}

\section{Results}

\subsection{ATSEmapper defines splicing events consistently across cells and datasets}
\todo{Fig.~\ref{fig:1}.}

\begin{figure*}[!t]
\centering
%\includegraphics[width=\textwidth]{Figures/fig1.pdf}
\caption{\textbf{ATSEmapper and the ATSE catalog.} \todo{caption from paper/captions/fig1.md}}
\label{fig:1}
\end{figure*}

\subsection{LeafletFA's per-cell representation reads junction usage and fits held-out reads}
% Held-out gain per read over the population null: \HeldoutGainPopulationTMS{} (TMS), \HeldoutGainPopulationAllen{} (Allen).
% Ratios redrawn: kNN drops by \KnnDropRatiosTMS{}. Seed range \SeedKnnRangeTMS{} / \SeedKnnRangeAllen{}.
% Per-cell imputed usage (rho = phi . psi) reproduces across 15 seeds: r \PerCellPsiSeedR{} over 105 fit pairs (Fig. 2e),
% while only 4 of the live-program clusters recur (check 8): programs are not identifiable, per-cell junction usage is.
\todo{Fig.~\ref{fig:2}.}

\begin{figure*}[!t]
\centering
%\includegraphics[width=\textwidth]{Figures/fig2.pdf}
\caption{\textbf{The model and what its per-cell representation reads.} \todo{caption; panel e: per-cell usage reproduces across seeds}}
\label{fig:2}
\end{figure*}

\subsection{On data simulated from the model, one fit recovers the programs at full-length depth}
% \SimPsiR{}, \SimCenteredCos{}, \SimProgramsRecovered{} programs, per-cell phi r \SimPhiR{}; \SimProgramsRecoveredEasySci{} at EasySci-like depth.
\todo{Fig.~\ref{fig:3}.}

\begin{figure*}[!t]
\centering
%\includegraphics[width=\textwidth]{Figures/fig3.pdf}
\caption{\textbf{Recovery on data simulated from the model.} \todo{caption}}
\label{fig:3}
\end{figure*}

\subsection{Programs that recur across seeds are broad lineages; imputation follows reads where a cell has them}
% Consensus over 15 seeds: \ConsensusLivePrograms{} live programs, \ConsensusRecurring{} recur, \ConsensusPassing{} keep their junctions;
% pseudobulk confirms \PseudobulkConfirmedBrain{} (brain) and \PseudobulkConfirmedEpithelial{} (epithelial).
% Fgfr2: covered mesenchymal \FgfrCoveredMesenchymal{}, uncovered \FgfrUncoveredMesenchymal{}.
\todo{Fig.~\ref{fig:4}.}

\begin{figure*}[!t]
\centering
%\includegraphics[width=\textwidth]{Figures/fig4.pdf}
\caption{\textbf{Programs that recur across seeds, and imputation at a known event.} \todo{caption}}
\label{fig:4}
\end{figure*}

\section{Discussion}
% Limits stated: duplicated reads cap per-cell resolution; programs beyond the broad lineages are seed-specific; depth limit (EasySci);
% the per-cell representation and the catalog are the supported outputs. Companion: SpliceVI \cite{SpliceVI2026}.
% One line (Karin, 2026-11-01): LeafletFA's posterior effect-size test for differential splicing is very conservative at this
% depth (no calls on a real-data null, few calls between cell types, not reproducible across seeds); DS is not part of this paper.
\todo{Discussion.}

\section{Data availability}
\todo{TMS (AWS), Allen (GEO GSE185862), catalog + fits on Zenodo.}

\section{Acknowledgements}
\section{Funding}

\bibliographystyle{oup-abbrvnat}
\bibliography{references}

\end{document}
